\chapter{Pancreatitis}

\medskip
\noindent\textit{\textbf{Under review.} This chapter is an AI-assisted educational draft; it is not a substitute for professional medical advice.}

\section*{Definition and Overview}
inflammation of the pancreas, most often related to gallstones or alcohol but also to medicines, high triglycerides, or other causes.

\section*{Clinical Features}
sudden severe upper-abdominal pain often radiating to the back, nausea, vomiting, fever, and abdominal tenderness.

\section*{When to Seek Medical Care}
Seek medical review for new, persistent, worsening, or function-limiting symptoms. Call 000 for severe breathing difficulty, collapse, sudden neurological change, severe chest or abdominal pain, or any rapidly worsening illness.

\section*{Diagnosis and Investigations}
clinical assessment, lipase, blood tests, ultrasound or CT when indicated, and investigation for gallstones or other causes.

\section*{Management}
hospital treatment may include fluids, analgesia, nutrition support, treatment of the cause, and management of complications; severe disease may require intensive care.

\section*{Complications}
Complications depend on the underlying cause and may include progression, effects on other organs, treatment adverse effects, or reduced quality of life. Early assessment can reduce preventable harm.

\section*{Prognosis and Self-Care}
Outcome depends on the cause, severity, complications, and response to treatment. Follow the care plan, keep appointments, avoid known triggers, and seek help if symptoms worsen.

\section*{Expanded clinical context}
Pancreatitis is a clinical condition or body-system topic. Interpretation should account for age, baseline health, medicines, comorbidities, goals, and access to care. This chapter supports discussion with a qualified clinician and does not establish a diagnosis.

\section*{Assessment and decision points}
Assess onset, progression, severity, functional impact, associated symptoms, past and family history, medicines, comorbidities, and relevant exposures. Document the main concern, time course, functional impact, relevant risks, and red flags. Use targeted investigations when their result is likely to change management.

\section*{Management and follow-up}
Management may include education, observation, lifestyle measures, medicines, rehabilitation, procedures, or referral, with benefits and risks discussed. Explain expected improvement, when reassessment is needed, and how follow-up can be accessed. Avoid starting, stopping, or sharing prescription medicines without professional advice.

\section*{Safety-netting}
Seek urgent care for sudden severe symptoms, chest pain, breathing difficulty, collapse, new neurological deficit, or rapidly worsening function.

\section*{Practical prevention}
Follow the care plan, attend monitoring, and arrange review for persistent or changing symptoms. Keep written instructions and seek review if the topic recurs, changes, or does not improve as expected.

\section*{Limitations}
This educational expansion should be checked against current local guidance and authoritative topic-specific sources.
\clearpage
